%%%%%
%
\section{Maximum-Echantillonnage d'un param\`etre {\tt float}}

\itemstyle{Op\'erateur~:} {\sf MaxEch}
\index{MaxEch}

%%%%%OBJET%%%%
\itemstyle{Objet~:}
Calcul des valeurs maxi par \'echantillon d'un
param\`etre \`a une fr\'equence donn\'ee, exemple
recherche des valeurs maxi toutes les nn minutes
pour une vitesse de vent.
Le temps en sortie est le temps de la valeur maxi
trouv\'e dans l'\'echantillon

%%%%%ARGUMENTS%%%%
\bigskip
%%%%%ENTREES%%%%
\begin{tabarg}
\multicolumn{3}{l}{\titretab {Entr\'ees}}\\
\hline
$X$ & Vecteur, {\tt float} & Param\`etre d'entr\'ee\\
$Time_{in}$ & Vecteur, {\tt double} & Vecteur temps correspondant au vecteur d'entr\'ee\\
$R_f$ & Coeff., {\tt double} & Rapport des fr\'equences ($ F_{sortie}\over
F_{entree} $) \\
$Decal$ & Coeff., {\tt int} & D\'ecalage (\textit{nombre de points})\\
\end{tabarg}

\medskip
%%%%%SORTIES%%%%
\begin{tabarg}
\multicolumn{3}{l}{\titretab {Sorties}}\\
\hline
$Time_{out}$ & Vecteur, {\tt double} & Vecteur temps correspondant aux valeurs maxi\\
$X_m$ & Vecteur, {\tt float} & R\'esultat (m\^eme unit\'e que $X$)\\
\end{tabarg}


%%%%%FORMULE%%%%
\itemstyle{Formule}

%%Max
Recherche de la valeur max dans chaque �chantillon
\[
X_m[i] = \max_{i/R_f \le j< {(i+1)}/{R_f}} (X[j]) %%\qquad j=\frac{i}{R_f}, \frac{i+1}{R_f}
%%i=0,Nbre\_Echant(X). R_f
\]
pour $i$ variant de 0 � $Nbre\_Echant(X). R_f$.


%%%%COMMENTAIRE%%%%
\itemstyle{Commentaire~:}

$R_f$ doit \^etre compris dans l'intervalle $]0, 1]$,
et \^etre l'inverse d'un entier. \newline
Le nombre d'\'echantillons du vecteur de sortie est le produit du
nombre d'\'echantillons du vecteur d'entr\'ee par $R_f$, arrondi
\`a l'entier sup\'erieur.

%%

%%%%%AUTEUR%%%%
\itemstyle{Auteur~:}
CNRM/GMEI/TRAMM

%%%%%REFERENCE%%%%
\itemstyle{R\'ef\'erences~:}
